% =====================================================================
% WitnessRAG: formulação teórica (Planejador, Executor, Refletor)
% Compilar: latexmk -pdf witnessrag-teoria.tex
% Descreve o controlador de prova como está implementado em 25/09/2026
% (wrag/methods/witnessrag.py::_retrieve_proof e módulos wrag/witness/*).
% =====================================================================
\documentclass[10pt,twocolumn]{article}
\usepackage[a4paper,top=1.3cm,bottom=1.5cm,left=1.4cm,right=1.4cm,columnsep=0.55cm]{geometry}
\usepackage[T1]{fontenc}
\usepackage[utf8]{inputenc}
\usepackage{lmodern}
\usepackage[brazilian]{babel}
\usepackage{amsmath,amssymb,amsthm,mathtools}
\usepackage{microtype}
\usepackage{booktabs,tabularx,array}
\usepackage{enumitem}
\usepackage{xcolor}
\usepackage{tikz}
\usetikzlibrary{arrows.meta,positioning,calc,fit,backgrounds,shapes.geometric}
\usepackage[ruled,vlined,linesnumbered]{algorithm2e}
\usepackage[hidelinks]{hyperref}
\usepackage{titlesec}
\usepackage{caption}
\captionsetup{font=small,labelfont=bf,skip=3pt}

\titleformat{\section}{\normalfont\large\bfseries}{\thesection}{0.5em}{}
\titlespacing*{\section}{0pt}{7pt plus 2pt}{3pt}
\titleformat{\subsection}{\normalfont\normalsize\bfseries}{\thesubsection}{0.5em}{}
\titlespacing*{\subsection}{0pt}{5pt plus 1pt}{2pt}
\titlespacing*{\paragraph}{0pt}{3pt plus 1pt}{0.6em}
\setlength{\parskip}{1.5pt}
\setlist{leftmargin=*,itemsep=0.5pt,topsep=2pt,parsep=0pt}

\theoremstyle{definition}
\newtheorem{definicao}{Definição}
\newtheorem{proposicao}{Proposição}
\newcommand{\rel}[1]{\textsf{#1}}
\newcommand{\V}[1]{\mathord{?}\mathit{#1}}
\newcommand{\op}[1]{\textsc{#1}}
\newcommand{\M}{\mathcal{M}}
\newcommand{\G}{\mathcal{G}}
\newcommand{\Pl}{\mathcal{P}}
\newcommand{\E}{\mathcal{E}}
\newcommand{\sem}{\operatorname{sim}}
\newcommand{\prox}{\operatorname{prox}}
\newcommand{\sq}{\hspace{0.08em}}
\DeclareMathOperator{\sig}{sig}
\DeclareMathOperator{\cost}{custo}
\DeclareMathOperator{\supp}{sup}
\newcolumntype{Y}{>{\raggedright\arraybackslash}X}

\SetAlCapFnt{\small}\SetAlCapNameFnt{\small}
\SetKwInOut{Input}{Entrada}\SetKwInOut{Output}{Saída}
\SetKw{Return}{devolva}\SetKw{Break}{pare}

\title{\vspace{-8mm}\textbf{WitnessRAG: memória de fatos datados e busca de testemunhas}\\[2pt]
\large Formulação com Planejador, Executor e Refletor}
\author{}
\date{\vspace{-10mm}}

\raggedbottom
\begin{document}
\maketitle


\begin{abstract}\noindent
O WitnessRAG guarda a conversa de um agente em duas camadas ligadas: texto
datado, que o leitor lê, e um grafo de fatos datados, em que a resposta é
\emph{provada}. Para cada pergunta, um \op{Planejador} escreve um plano de
busca, um \op{Executor} percorre o grafo seguindo o plano até achar uma
\emph{testemunha} (os poucos fatos dos quais a resposta decorre) e um
\op{Refletor} confere a prova no texto original ou explica ao Planejador por
que ela falhou. O contexto do leitor só muda quando há prova confirmada.
\end{abstract}

\begin{figure*}[t]
\centering
\begin{tikzpicture}[
  ent/.style={draw,ellipse,thick,fill=orange!12,font=\small,inner sep=1.5pt,minimum height=6mm},
  e/.style={-{Stealth[length=2mm]},thick},
  eA/.style={e,draw=blue!70!black,line width=1.3pt},
  eC/.style={e,draw=red!70!black,line width=1.3pt},
  el/.style={font=\scriptsize,fill=white,inner sep=1pt,align=center},
  pas/.style={draw,rounded corners=2pt,fill=gray!8,font=\scriptsize,align=left,inner sep=3pt,text width=4.3cm}]
\node[ent] (bet) at (-7.0,0.9) {Beth};
\node[ent] (car) at (-3.4,0.9) {Caroline};
\node[ent] (bos) at (-7.0,-1.1) {Boston};
\node[ent] (mel) at (-0.6,0) {Melanie};
\node[ent] (pot) at (-4.0,-1.25) {pottery class};
\node[ent] (mon) at (3.1,1.35) {mountains};
\node[ent] (bea) at (3.4,0) {beach};
\node[ent] (for) at (3.1,-1.35) {forest};
\draw[eC] (bet)--node[el,above=1pt,fill=none]{$f_4$ mentor of, 8/5}(car);
\draw[eC] (bet)--node[el,left=1pt,fill=none]{$f_5$ move to\\ jul/2023}(bos);
\draw[e] (mel)--node[el,above,sloped,pos=0.55,fill=none]{$f_6$ sign up for, 2/7}(pot);
\draw[eA] (mel)--node[el,above,sloped,pos=0.38,fill=none]{$f_1$ camp at, 19--25/6}(mon);
\draw[eA] (mel)--node[el,above,fill=none]{$f_2$ camp at, 6/7}(bea);
\draw[eA] (mel)--node[el,below,sloped,fill=none]{$f_3$ camp at, 15/7}(for);
\node[pas,anchor=west] (p) at (4.4,0.0) {%
\textbf{origem $\pi(f)$: trecho, fala, sessão}\\[1pt]
$f_1$: $p_2$, $u_{4:6}$, 27/6: ``I took my fam camping in the mountains \emph{last week}''\\[1pt]
$f_2$: $p_3$, $u_{6:16}$, 6/7: ``\dots my family camping at the beach''\\[1pt]
$f_3$: $p_4$, $u_{8:3}$, 15/7: ``we went camping in the forest''\\[1pt]
$f_4$: $p_1$, 8/5;\quad $f_5$: $p_5$, 14/8 (``\emph{last month}'');\quad $f_6$: $p_3$, 2/7};
\end{tikzpicture}
\caption{Memória didática no formato do LoCoMo. Nós são classes de identidade; arestas são fatos com relação e data do
\emph{evento} ($f_1$: ``last week'' resolvida pela sessão de 27/6; $f_5$: ``last month'' dito em 14/8). Em azul, as
testemunhas do Exemplo A; em vermelho, a do Exemplo C. À direita, o ponteiro $\pi(f)$ para a camada de texto.}
\label{fig:grafo}
\vspace{-3mm}
\end{figure*}

% =====================================================================
\section{A ideia}
\label{sec:ideia}

Um agente que conversa durante meses não consegue reler tudo a cada pergunta:
ele escolhe poucos trechos ($k=5$) e um \emph{leitor} (o LLM) responde a partir
deles. A escolha usual é por similaridade com a pergunta, e ela falha de três
formas: quando a resposta está \emph{espalhada} (três acampamentos em três
sessões), quando exige \emph{encadear} fatos (a cidade da mentora de Caroline
está num trecho que nem cita Caroline) e quando depende de \emph{datas}
(``em junho'' contra ``em julho''). Nos três casos falta a mesma coisa: saber
se os trechos escolhidos \emph{bastam} para responder.

O WitnessRAG troca a pergunta ``o que se parece com a pergunta?'' por ``que
fatos provam a resposta?''. O conjunto mínimo de fatos que prova uma resposta
é a sua \emph{testemunha}, termo da teoria de proveniência em bancos de dados.
Ela serve para duas coisas: indica \emph{quando parar} de procurar (achou a
prova) e \emph{o que o contexto precisa conter} (os trechos de onde vêm esses
fatos). O processo imita um investigador: planeja o que procurar, procura, e
confere o que achou antes de concluir; se falhar, entende o motivo e tenta de
novo.

% =====================================================================
\section{A memória: o que é guardado}
\label{sec:memoria}

A memória é construída uma vez por conversa, antes de qualquer pergunta. Ela
tem duas camadas: o \emph{texto}, para ser lido, e o \emph{grafo}, para
raciocinar. Seja $\mathbb{D}$ o conjunto de intervalos de dias $[d_1,d_2]$
(um dia, uma semana, um mês).

\begin{definicao}[Fala e trecho]
Uma \emph{fala} é $u=(\mathrm{id},\,\mathrm{falante},\,\mathrm{texto},\,d_u)$,
com $d_u$ a data da sessão. Um \emph{trecho} $p$ é uma sequência contígua de
falas com até 2048 tokens, com intervalo
$\tau(p)=[\min_{u\in p}d_u,\max_{u\in p}d_u]$ e importância $\iota(p)$ (média
de $\iota(u)$ nas falas, normalizada na memória).
\end{definicao}

\emph{Importância} mede o quanto o falante marcou um evento como significativo.
Usamos a intensidade emocional do texto, contada por um léxico pequeno e
auditável: $E(u)$ soma os pesos das palavras emotivas (1 ou 1{,}5), dos
intensificadores (0{,}3) e das exclamações (0{,}5, até 3), e
\begin{equation}\label{eq:imp}
\iota(u)=1-e^{-E(u)/\kappa}\in[0,1),\qquad \kappa=2.
\end{equation}
A exponencial satura: a primeira palavra emotiva pesa muito, a décima quase
nada. ``Foi incrível, estou tão feliz!'' fica perto de 0{,}9; ``fui ao
mercado'', em 0.

\begin{definicao}[Fato]
Um \emph{fato} é
$f=(s_f,\,r_f,\,o_f,\,\tau(f),\,\iota(f),\,c_f,\,\pi(f))$: sujeito e objeto,
relação em texto livre, intervalo do \emph{evento} $\tau(f)\in\mathbb{D}$,
importância $\iota(f)=\iota(u_f)$, confiança da extração $c_f\in[0,1]$ e origem
$\pi(f)=(p_f,u_f)$, o trecho e a fala de onde veio.
\end{definicao}

Em palavras: um fato é uma frase curta do tipo ``Melanie acampou nas
montanhas'', acompanhada de \emph{quando} aconteceu, \emph{quão marcante} foi
e \emph{de onde} saiu. Os fatos são extraídos por um LLM em janelas curtas do
diálogo, usando o falante como sujeito. A relação é o verbo que o texto usou
(``camp at'', ``move to''); não existe lista fixa de relações. A data é a do
\emph{evento}, não a da conversa: a expressão de tempo é resolvida pela data
da fala, $\tau(f)=\mathrm{res}(\epsilon_f,d_{u_f})$, e ``last week'' dito em
27/6/2023 vira $[19/6,25/6]$. Sem expressão de tempo, o fato recebe o dia da
sessão. A fala de origem $u_f$ é a fala do trecho que mais compartilha
palavras com o fato.

\begin{definicao}[Grafo de memória]
$\G=(V,F,\sigma,\omega)$ é um multigrafo dirigido. Os nós $V$ são
\emph{classes de identidade} de nomes; cada fato $f\in F$ é uma aresta de
$\sigma(f)=[s_f]$ para $\omega(f)=[o_f]$, rotulada por $r_f$ e com os
atributos $(\tau,\iota,c,\pi)$.
\end{definicao}

Pessoas e coisas viram nós; cada fato vira uma seta do sujeito para o objeto
(Fig.~\ref{fig:grafo}). É um \emph{multi}grafo porque o mesmo par pode ter
várias setas: Melanie acampou em lugares e datas diferentes, e cada relato
guarda sua data e sua origem. O ponteiro $\pi(f)$ é a ponte entre as camadas:
a prova é encontrada no grafo, mas o leitor recebe o texto original.

\paragraph{Quando dois nomes são a mesma coisa.} ``Mel'' e ``Melanie'' devem
ser um só nó, mas ``primeira metade do jogo'' e ``segunda metade do jogo'' não,
embora sejam quase idênticos para um embedding. Por isso dois nomes só se
unem (union--find) se $\cos(\phi(n),\phi(n'))\ge0{,}80$ \emph{e} as palavras
de um estão contidas nas do outro. Unir errado fabricaria provas falsas;
deixar de unir só perde uma prova. Flexões de uma relação (``paint'',
``painted'') compartilham uma família, e relações diferentes são comparadas
por significado:
$\sem_R(r,r')=1$ se são da mesma família e $\cos(\phi(r),\phi(r'))$ caso
contrário, com $\phi$ o embedding (BGE).

\begin{table}[t]
\centering\small
\setlength{\tabcolsep}{3pt}
\begin{tabularx}{\columnwidth}{@{}l Y Y@{}}
\toprule
Componente & Guarda & Serve para \\
\midrule
Fala $u$ & id, falante, texto, $d_u$, $\iota(u)$ & ser a prova legível (verificação, bloco de falas) \\
Trecho $p$ & falas, $\tau(p)$, $\iota(p)$, $\phi(p)$ & ser lido; busca densa+BM25 \\
Nó $[n]$ & nome canônico, variantes, $\phi(n)$ & ancorar constantes do plano \\
Aresta $f$ & $s,r,o,\tau,\iota,c,\pi$, $\phi(f)$ & compor a prova (junção) \\
\bottomrule
\end{tabularx}
\caption{O que cada componente da memória $\M=(U,P,\G)$ guarda. O ponteiro
$\pi(f)$ liga a camada lógica à camada de texto.}
\label{tab:mem}
\vspace{-3mm}
\end{table}

% =====================================================================
\section{Planos, testemunhas e pontuação}
\label{sec:plano}

\begin{definicao}[Plano]
Um \emph{plano} é $\Pl=(Q,\alpha,\theta,I^\star,\mathbf w)$, com
\begin{equation*}
Q(\V x)\leftarrow \textstyle\bigwedge_{i=1}^{n} r_i(t_i,t'_i)\,[@\,\V t],\qquad n\le4.
\end{equation*}
\end{definicao}

O plano é a pergunta reescrita na língua do grafo. $Q$ é uma lista de fatos a
encontrar, ligados por ``e'' ($\wedge$); termos com ``?'' são incógnitas e os
demais são nomes. ``Onde mora a irmã de Ana?'' vira
$\rel{sister\,of}(\V y,\text{Ana})\wedge\rel{live\,in}(\V y,\V x)$: a
incógnita $\V y$ aparece nos dois átomos e \emph{amarra} a cadeia (a irmã do
primeiro fato tem de ser a mesma pessoa do segundo). O marcador $@\,\V t$ pede
a data do fato, para perguntas ``quando''. Os demais campos dizem que
\emph{forma} tem a resposta e \emph{onde procurar}: $\alpha$ é a agregação
(um valor, um conjunto, uma contagem), $\theta(\V x)$ um tipo opcional
(``arte marcial''), $I^\star$ o período de referência (\emph{hoje},
\emph{início} ou uma \emph{janela} citada, como ``junho de 2023'') e
$\mathbf w$ os pesos da pontuação abaixo.

\begin{definicao}[Testemunha]
Uma atribuição $\beta$ das incógnitas e um conjunto de fatos
$W=\{f_1,\dots,f_n\}$ formam uma \emph{testemunha} da resposta
$a=\beta(\V x)$ se cada átomo $i$ é realizado por $f_i$
($\sigma(f_i)=\beta(t_i)$, $\omega(f_i)=\beta(t'_i)$,
$r_{f_i}\approx r_i$ e, havendo $\V t$, $\tau(f_i)=\beta(\V t)$) e nenhuma
parte própria de $W$ basta (mínima).
\end{definicao}

Em palavras: uma testemunha é uma forma de preencher todas as incógnitas com
nós do grafo tal que cada átomo do plano corresponda a uma seta que existe de
fato. Ela é, literalmente, a prova. Uma resposta pode ter várias
testemunhas (várias provas independentes); o conjunto de respostas é
$\mathrm{Ans}(Q,\G)=\{\beta(\V x):\exists(\beta,W)\}$.

\paragraph{Uma pontuação para tudo.} Toda escolha de ordem, de trechos e de
fatos, usa a mesma fórmula, cujos pesos vêm do plano:
\begin{equation}\label{eq:score}
S_{\Pl}(x)=w_{\text{sim}}\sem(x)+w_{\text{imp}}\,\iota(x)+w_{\text{t}}\prox(x),
\end{equation}
\begin{equation}\label{eq:prox}
\prox(x)=\exp\!\big(-\mathrm{dist}(\tau(x),I^\star)/\lambda\big).
\end{equation}
Os três sinais respondem a perguntas simples: \emph{parece com o que procuro?}
($\sem$: para um trecho, busca densa+BM25 fundida por RRF; para um fato, o
casamento com o átomo, Eq.~\ref{eq:match}); \emph{foi marcante?} ($\iota$);
\emph{aconteceu perto do período pedido?} ($\prox$). A proximidade vale 1
dentro do período e cai exponencialmente com os dias de distância, com escala
$\lambda=\max(7,|I^\star|/2)$ para uma janela e $\max(7,0{,}25\,|\M|)$ para
\emph{hoje} ou \emph{início}. A mesma fórmula cobre três ideias: recência
(período \emph{hoje}), ``a primeira vez'' (\emph{início}) e ``em junho''
(\emph{janela}). Os pesos têm três níveis: tempo forte vale 0{,}3,
importância forte 0{,}1, normal vale 0; eles somam 1 e a similaridade fica
sempre com pelo menos metade, para que nada irrelevante passe à frente só por
ser recente ou emotivo. No nível normal, a ordem é exatamente a da busca
híbrida.

% =====================================================================
\section{A busca: Planejador, Executor e Refletor}
\label{sec:busca}

\begin{figure}[t]
\centering
\begin{tikzpicture}[
  box/.style={draw,rounded corners=2pt,thick,minimum height=7mm,minimum width=16mm,font=\sffamily\small,align=center},
  llm/.style={box,fill=orange!15}, det/.style={box,fill=blue!8},
  arr/.style={-{Stealth[length=1.8mm]},thick}, lab/.style={font=\scriptsize,align=center}]
\node[det] (ex0) {Executor\\[-1pt]\scriptsize busca $\Pl_0$};
\node[llm,right=9mm of ex0] (pl) {Planejador\\[-1pt]\scriptsize $\Pi$};
\node[det,right=9mm of pl] (ex) {Executor\\[-1pt]\scriptsize $\Xi$};
\node[llm,below=5mm of ex] (re) {Refletor\\[-1pt]\scriptsize $\rho$};
\node[det,below=4mm of re] (ct) {Contexto\\[-1pt]\scriptsize $C$};
\node[font=\small,left=8mm of ct] (rd) {leitor};
\draw[arr] (ex0)--node[lab,above]{$\E_0$}(pl);
\draw[arr] (pl)--node[lab,above]{$\Pl_j$}(ex);
\draw[arr] (ex)--node[lab,right]{$O_j$}(re);
\draw[arr] (re)--node[lab,right]{$W^\star$ confirmada}(ct);
\draw[arr] (ct)--(rd);
\draw[arr,dashed] (re.west)-|node[lab,pos=0.3,below]{feedback $\varphi_j$, lacuna}node[lab,pos=0.8,left]{até\\$R$ ciclos}(pl.south);
\end{tikzpicture}
\caption{Laço de busca. Laranja chama o LLM; azul é determinístico.}
\label{fig:loop}
\vspace{-3mm}
\end{figure}

A busca é um diálogo entre três papéis (Fig.~\ref{fig:loop},
Alg.~\ref{alg:loop}). O Planejador e o Refletor usam o LLM; o Executor é
determinístico, só índices e contas. Antes do laço, uma busca por
similaridade pura (plano nulo $\Pl_0$, $\mathbf w=(1,0,0)$) traz as
\emph{evidências} $\E_0$: os $N=20$ trechos mais parecidos com a pergunta. O
laço roda no máximo $R=2$ ciclos.

\paragraph{Planejador $\Pi$: decide o que procurar.}
$\Pl_j=\Pi(q,\;\mathrm{Fatos}(\E),\;\mathcal V,\;\varphi_{<j})$. Ele não
planeja às cegas: vê até 30 fatos das evidências (com datas), o vocabulário de
relações e nomes $\mathcal V$ que a memória realmente usa e o feedback dos
ciclos anteriores. Assim escreve o plano ``com as palavras da memória''. Nunca
vê categoria da pergunta nem resposta.

\paragraph{Executor $\Xi$: segue o plano no grafo.} Em quatro passos.

\emph{(1) Ancorar os nomes.} Cada nome do plano é ligado aos nós
compatíveis: o nó idêntico e até 5 nós com cosseno $\ge0{,}55$.

\emph{(2) Achar candidatos para cada átomo.} Entre as setas que saem ou chegam
às âncoras, cada fato $f$ recebe uma nota de casamento com o átomo $a=r(t,t')$:
\begin{equation}\label{eq:match}
m(a,f)=\frac{0{,}45\,\sem_R(r,r_f)+0{,}35\,\bar s_{\text{arg}}+0{,}20\,\sem_V}{0{,}45+0{,}35\,\mathbb 1_{\text{const}}+0{,}20}.
\end{equation}
Pesam a relação ($\sem_R$), os nomes ($\bar s_{\text{arg}}$, similaridade das
âncoras) e a frase inteira ($\sem_V$, cosseno entre ``Melanie camp at ?x'' e
o fato verbalizado). É uma média ponderada, e o denominador só a renormaliza
quando o átomo não tem nome. Há dois vetos: relação com $\sem_R<0{,}65$, ou
nome no lado errado da seta, dão $m=0$. Os 60 melhores por
$s(a,f)=w_{\text{sim}}m(a,f)+w_{\text{imp}}\iota(f)+w_{\text{t}}\prox(f)$
formam $\Gamma(a)$.

\emph{(3) Juntar os átomos.} A busca anda camada por camada, começando pelo
átomo com mais nomes (o mais restrito). Cada estado parcial $(\beta,W)$ só
aceita um novo fato se as incógnitas compartilhadas concordam: se $\V y$ já é
Beth, o próximo fato tem de sair de Beth. Os estados são ordenados por
\begin{equation}\label{eq:cost}
\cost(W)=-\textstyle\sum_{f\in W}\log s(a_f,f)+0{,}15\,\big|\{p_f:f\in W\}\big|,
\end{equation}
que prefere fatos bem casados e provas espalhadas por poucos trechos (cada
trecho distinto custa 0{,}15), com feixe de 400 estados. Se uma camada fica
vazia, a busca para e registra a \emph{lacuna} $g=(a_i,\beta)$: qual átomo
faltou e o que já estava ligado. Saber \emph{onde} a prova quebrou é o que
permite corrigir o plano.

\emph{(4) Testar se a prova é boa o bastante.} Cada testemunha recebe o
suporte $\supp(W)=\prod_{f\in W}m(a_f,f)\,c_f$; o produto faz a prova valer o
seu elo mais fraco. Um plano de \emph{valor} só é aceito se é seletivo: até 3
respostas e 3 testemunhas, sem cortes, $\max\supp\ge0{,}45$ e trechos novos
dentro do orçamento $k_W$. Muitas respostas indicam um plano vago, não uma
prova. Num plano de \emph{conjunto}, muitas respostas são esperadas, então
cada item precisa da própria testemunha com $\supp\ge0{,}45$ (com tipo, os 12
melhores seguem; sem tipo e com mais de 12 itens, o plano é vago).

\paragraph{Refletor $\rho$: confere e explica.} Tem dois papéis.
\emph{Verificar}: se a prova vai trazer trecho novo ao contexto, o LLM lê a
pergunta, o plano e, para cada resposta, os fatos com suas datas \emph{e as
falas originais}; aceita ou recusa item a item (entidade, período ou tipo
errados, ou falta de apoio). Ele julga, não procura; na dúvida, aceita. Se a
prova já está no contexto, não há o que verificar. \emph{Diagnosticar}: se
não houve prova, ou ela foi recusada, transforma o motivo numa instrução
para o próximo plano: ``nenhum fato casou \rel{live in} para Beth'' (lacuna),
``8 respostas: seja mais específico'', ``só casamentos fracos: use as
relações da memória'', ``recusada: período errado''. A lacuna vira também uma
sonda textual (``Beth live in'') que traz novos fatos às evidências. Plano
repetido encerra o laço.

\begin{algorithm}[t]
\small
\caption{Laço de prova (uma pergunta $q$)}\label{alg:loop}
\Input{$q$, memória $\M$, $k$, $R$}
$\E\gets\operatorname{top}_N S_{\Pl_0}$,\ \ $\varphi\gets[\,]$\;
\For{$j=1,\dots,R$}{
  $\Pl_j\gets\Pi(q,\mathrm{Fatos}(\E),\mathcal V,\varphi)$\;
  \lIf{$\Pl_j$ inválido}{$\varphi\mathrel{+}=$erro; \textbf{continue}}
  \lIf{$\Pl_j$ repetido}{\Break}
  $\E\gets\E\cup\operatorname{top}_N S_{\Pl_j}$,\ \ $O_j\gets\Xi(\Pl_j,\E)$\;
  \If{$O_j$ seletivo}{
    $W^\star\gets\rho_{\text{verif}}(q,\Pl_j,O_j)$\;
    \lIf{$W^\star\neq\emptyset$}{\Return $C(\Pl_j,W^\star)$}}
  $\varphi\mathrel{+}=\rho_{\text{diag}}(O_j)$ \tcp*{causa e sonda}
}
\Return $\operatorname{top}_k S_{\Pl}$ \tcp*{sem prova: busca pura}
\end{algorithm}

\paragraph{Montar o contexto.} O leitor recebe $B=\operatorname{top}_k S_\Pl$
($k=5$) e, havendo prova confirmada, os trechos da testemunha que faltam
entram \emph{pela cauda}, até $k_W$ trocas ($k_W=2$ para cadeia, interseção
ou conjunto; $1$ para um fato e um valor). Os $\max(1,k-k_W)$ primeiros
trechos nunca mudam. Provas que não cabem entram como um bloco curto das
falas de origem (até 2400 caracteres), sem triplas nem respostas extraídas,
para o leitor não copiar a extração. Perguntas hipotéticas (``Would X\dots?'')
recebem até 8 falas da vizinhança de X, como premissas e não como prova.

\begin{proposicao}[Conservadorismo]
Sem prova confirmada e com pesos normais, os $k$ trechos de $C$ são os do
híbrido; em qualquer caso, $|C\setminus B|\le k_W$.
Ou seja, o método só \emph{acrescenta} evidência provada.
\end{proposicao}

% =====================================================================
\section{Exemplos de percurso}
\label{sec:exemplos}

Os exemplos usam a memória da Fig.~\ref{fig:grafo}. Limiares e pesos são os
do código; as similaridades são ilustrativas; $c_f=0{,}9$.

\paragraph{A: um salto, resposta em conjunto.} \emph{``Onde Melanie já
acampou?''} Plano: $\rel{camp\,at}(\text{Melanie},\V x)$, conjunto, tipo
``place''. O Executor ancora Melanie e olha as quatro setas que saem dela.
Para $f_2$, relação idêntica, nome idêntico e frase parecida (0{,}82):
$m=0{,}45+0{,}35+0{,}20\cdot0{,}82=0{,}964$. A seta $f_6$ (\rel{sign up for})
é vetada, pois a relação não se parece com \rel{camp at}. Sobram três
testemunhas de um fato só, montanhas, praia e floresta, cada uma com
$\supp=0{,}964\cdot0{,}9\approx0{,}87$. O Refletor lê as três falas e confirma.
Se o trecho $p_4$ (floresta) não estava entre os 5, ele entra no lugar do
último.

\paragraph{B: o tempo decide.} \emph{``Quando Melanie acampou em junho?''}
Plano: $\rel{camp\,at}(\text{Melanie},\V x)\,@\,\V t$, resposta $\V t$,
janela junho ($\lambda=15$ dias), tempo forte, $\mathbf w=(0{,}7;0;0{,}3)$.
Os três fatos casam igualmente bem ($m\approx0{,}96$); só a data os separa:
\begin{center}\small
\begin{tabular}{@{}lccc@{}}
\toprule
 & dias fora & $\prox$ & $s=0{,}7m+0{,}3\prox$ \\
\midrule
$f_1$ (19--25/6) & 0 & 1{,}000 & \textbf{0{,}972} \\
$f_2$ (6/7) & 6 & 0{,}670 & 0{,}873 \\
$f_3$ (15/7) & 15 & 0{,}368 & 0{,}782 \\
\bottomrule
\end{tabular}
\end{center}
A resposta é ``19--25 June 2023'', a semana antes de 27/6. Note que a data
correta existe graças à construção da memória: a fala dizia apenas ``last week''.

\paragraph{C: uma cadeia que precisa de reflexão.} \emph{``Para que cidade
se mudou a mentora de Caroline?''}\\
\emph{Ciclo 1.} Plano:
$\rel{mentor\,of}(\V y,\text{Caroline})\wedge\rel{live\,in}(\V y,\V x)$. O
Executor começa pelo átomo com nome: entre as setas que chegam a Caroline,
$f_4$ casa ($m=0{,}98$) e fixa $\V y=$ Beth. Na segunda camada, a única seta
que sai de Beth é $f_5$ (\rel{move to}), mas
$\sem_R(\rel{live in},\rel{move to})=0{,}58<0{,}65$: veto, camada vazia.
Lacuna: $(\rel{live in},\ \V y{=}\text{Beth})$. O Refletor diz ao Planejador
``nenhum fato casou \rel{live in} para Beth'', e a sonda ``Beth live in''
põe o fato (Beth | move to | Boston) entre as evidências mostradas.\\
\emph{Ciclo 2.} O Planejador reescreve com a relação da memória:\\[1pt]
{\scriptsize\ttfamily\setlength{\parskip}{0pt}%
\hspace*{1mm}"atoms":[\{"relation":"mentor of","subject":"?y",\\
\hspace*{11mm}"object":"Caroline"\},\\
\hspace*{11mm}\{"relation":"move to","subject":"?y","object":"?x"\}]}\\[2pt]
Agora $m(f_5)=(0{,}45\cdot1+0{,}20\cdot0{,}85)/0{,}65=0{,}954$ (sem nome no
átomo, a média é renormalizada). A testemunha é $W=\{f_4,f_5\}$, com
$\V y=$ Beth e $\V x=$ Boston, $\supp=0{,}98\cdot0{,}954\cdot0{,}9^2=0{,}76$ e
$\cost=0{,}02+0{,}05+0{,}15\cdot2=0{,}37$. O Refletor lê as duas falas e
confirma; o trecho $p_5$ entra na 5.ª posição. O leitor recebe um trecho que
a similaridade não traria: ``Boston'' nunca aparece perto de ``Caroline''.

\paragraph{Propriedades e limites.} O laço custa no máximo $2R=4$ chamadas de
LLM além do leitor, e nada nele lê categoria ou resposta do benchmark. Com
casamento exato e sem cortes, o Executor devolve exatamente
$\mathrm{Ans}(Q,\G)$; no modo semântico (padrão), é uma aproximação. A prova
herda os erros da extração, e um conjunto contém o que a memória prova, sem
garantir que nada ficou de fora.

\end{document}
